\section{Equações do Movimento de Atitude para uma Missão de RVD/B}
\label{sec:equacoes_movimento_atitude}

O problema de atitude em missões de RVD/B exige a descrição da orientação da espaçonave visitante em relação à espaçonave alvo.
Essa descrição deve ser consistente tanto no referencial ECI quanto LVLH.
No ECI, a atitude é descrita em relação ao centro da Terra, enquanto que no LVLH, a descrição da atitude se relaciona com as necessidades de alinhamento para aproximação, captura e acoplamento.
Nesta trabalho, foi adotado o método de Euler para a formulação das equações dinâmicas da atitude.
A presença de um manipulador robótico acoplado à base da espaçonave visitante imputa uma complexidade adicional, pois a movimentação de seus elos altera a distribuição de massa e a matriz de inércia do sistema.

\subsection{Equações do Movimento de Atitude da Espaçonave Visitante no Referencial Inercial}
\label{sec:atitude_inercial}

A formulação matemática de atitude da espaçonave visitante é realizada partir das equações de Euler aplicadas a um corpo rígido.
Essas equações são descritas em torno do centro de massa (\emph{Center of Mass} ou CoM) e expressas no sistema da própria espaçonave.
As equações diferenciais de movimento rotacional assumem a forma:

\begin{equation}
\mathbf{I}\,\dot{\vec{\omega}}
+ \vec{\omega}\times \big(\mathbf{I}\,\vec{\omega}\big)
= \vec{M},
\label{eq:euler_corpo_rigido}
\end{equation}

em que $\mathbf{I}$ é a matriz de inércia em relação ao centro de massa, $\vec{\omega}$ é o vetor de velocidade angular e $\vec{M}$ são os momentos aplicados ao corpo.

\subsubsection{Matriz de inércia variável}

No caso da espaçonave visitante, que possui um MR acoplado em sua base, a matriz de inércia equivalente não é constante, pois a movimentação dos elos modifica a distribuição de massa, sendo assim, a inércia total do sistema pode ser representada por:

\begin{equation}
\mathbf{I}(t) = \mathbf{I}_{base} + \mathbf{I}_{manip}(t),
\label{eq:inercia_total}
\end{equation}

em que $\mathbf{I}_{base}$ é a matriz associada à base da espaçonave e $\mathbf{I}_{manip}(t)$ é a contribuição variável do manipulador robótico.

\subsubsection{Dinâmica acoplada com o manipulador}

Os movimentos do MR geram torques que devem ser incluídos na equação dinâmica da atitude, \eqref{eq:euler_corpo_rigido}.
Deve-se considerar tanto os torques de controle externos quanto os torques internos de reação do MR, dada por:

\begin{equation}
\mathbf{I}(t)\,\dot{\vec{\omega}}
+ \vec{\omega}\times \big(\mathbf{I}(t)\,\vec{\omega}\big) 
= \vec{M}_{controle} + \vec{M}_{manip},
\label{eq:dinamica_acoplada}
\end{equation}

onde $\vec{M}_{controle}$ são os torques de controle aplicados à espaçonave e $\vec{M}_{manip}$ são os torques produzidos pelas juntas do MR durante sua movimentação. 

\subsection{Equações Diferenciais da Cinemática}
\label{sec:equacoes_cinematica}

A atitude da espaçonave visitante foi descrita utilizando os quaternions unitários por conta da ausência de singularidades (\emph{gimbal lock}) e maior eficiência computacional em simulações numéricas.

\subsubsection{Quaternions unitários}

Um quaternion unitário é definido como:

\begin{equation}
\label{eq:quaternion_unitario_definicao}
\mathbf{q} =
\begin{bmatrix}
q_0 \\ q_1 \\ q_2 \\ q_3
\end{bmatrix}
=
\begin{bmatrix}
\cos\frac{\alpha}{2} \\
u_x \sin\frac{\alpha}{2} \\
u_y \sin\frac{\alpha}{2} \\
u_z \sin\frac{\alpha}{2}
\end{bmatrix},
\end{equation}

onde $\alpha$ é o ângulo de rotação e $\vec{u}=(u_x,u_y,u_z)^T$ é o eixo unitário.
A condição de normalização é:

\begin{equation}
\label{eq:normalizacao_quaternion}
\lVert \mathbf{q} \rVert^2 = q_0^2+q_1^2+q_2^2+q_3^2 = 1.
\end{equation}

\subsubsection{Cinemática diferencial}

Considerando a velocidade angular $\vec{\omega}$ expressa no corpo da espaçonave, a cinemática do quaternion é dada por:

\begin{equation}
\label{eq:qdot}
\dot{\vec{q}} = \tfrac{1}{2}\,\Omega(\vec{q})\,\vec{\omega},
\end{equation}

sendo:

\begin{equation}
\Omega(\vec{q}) =
\begin{bmatrix}
-{\vec \epsilon\,}^T \\
q_0 \mathbf{I}_3 + [\vec \epsilon\,]_\times
\end{bmatrix},
\qquad
\vec{q} = \begin{bmatrix} q_0 \\
\vec{\epsilon} \end{bmatrix},
\end{equation}
onde $[\vec \epsilon\,]_\times$ é a matriz antissimétrica associada ao vetor $\vec \epsilon$.

\subsubsection{Acoplamento com a dinâmica}

No caso de corpo rígido, a equação \eqref{eq:qdot} descreve a evolução temporal da atitude a partir da velocidade angular. 
No entanto, quando o MR está em movimento.
Como visto em \eqref{eq:dinamica_acoplada}, a velocidade angular $\vec{\omega}$ não é independente; ela resulta da integração das equações de Euler com inércia variável.
A solução das equações diferenciais da cinemática depende da evolução conjunta de $\vec{\omega}$, a qual é influenciada pelas variáveis generalizadas do MR.
Esse acoplamento implica que o modelo da atitude da espaçonave visitante não é o de um corpo rígido, mas de um sistema dinâmico composto, em que a matriz inversa de inércia fornece as equações diferenciais completas da cinemática.

A equação de Euler modificada, apresentada em \eqref{eq:dinamica_acoplada}, pode ser reescrita isolando $\dot{\vec{\omega}}$:

\begin{equation}
\label{eq:euler_modificada}
\dot{\vec{\omega}}
= \mathbf{I}^{-1}(q)\Big(\vec{M}_{controle} + \vec{M}_{manip} - \vec{\omega} \times \big(\mathbf{I}(q)\vec{\omega}\big)\Big).
\end{equation}

onde $\mathbf{I}(q)$ é função dos ângulos e velocidades das juntas do manipulador robótico, de modo que tanto a matriz de inércia quanto sua inversa variam no tempo, sendo assim, a integração de \eqref{eq:qdot} e \eqref{eq:euler_modificada} fornece as equações diferenciais da cinemática da atitude da espaçonave visitante com a movimentação realizada pelo MR.

\subsection{Equações do Movimento da Atitude Relativa entre as espaçonaves visitante e alvo}

A atitude relativa entre duas espaçonaves descreve a orientação do \emph{chaser} em relação ao \emph{target}.
Essa relação fornece medidas de erro de atitude para que a lei de controle possa corrigir as orientações relativas.
Sejam $q_t$ e $q_c$ os quaternions unitários das espaçonaves \emph{target} e \emph{chaser}, ambos definidos em relação ao mesmo referencial inercial. 
A orientação relativa pode ser expressa por:

\begin{equation}
    \label{eq:quat_relativo}
q_{rel}
=
q_t^{-1} \otimes q_c,
\end{equation}

onde o termo $q_t^{-1}$ é o inverso (no caso dos quaternions unitários, conjugado) do quaternion do \emph{target}, $q_c$ é o quaternion do \emph{chaser}, e $\otimes$ é o produto de quaternions.
Ademais, a atitude relativa não corresponde à atitude absoluta do visitante nem a do alvo, mas à diferença entre ambas em relação ao ECI. 
Similarmente, ao movimento relativo de posição, definido pela subtração dos vetores $\vec r_c - \vec r_t$, a atitude relativa resulta da combinação de duas atitudes absolutas expressas em quaternions. 
A condição de sincronização em uma missão de RVD/B ocorre quando o quaternion relativo converge para a identidade unitária:

\begin{equation}
\label{eq:qrel_sincronizacao}
q_{rel} \;\;\longrightarrow\;\; 
\begin{bmatrix}1 \\ 0 \\ 0 \\ 0\end{bmatrix}.
\end{equation}

Isso significa que a orientação do visitante se tornou idêntica à orientação do alvo. 
Portanto, o erro de atitude é zero, e as espaçonaves estão alinhadas.
Sendo assim, o quaternion $q_{rel}$ representa a rotação necessária para alinhar as espaçonaves.
A parte escalar $q_{rel}$ é descrita por:

\begin{equation}
q_{rel} =
\begin{bmatrix}
\cos\!\left(\tfrac{\alpha}{2}\right) \\[4pt]
\vec{u}\,\sin\!\left(\tfrac{\alpha}{2}\right)
\end{bmatrix},
\end{equation}

onde $\alpha$ é o ângulo relativo entre os dois referenciais e $\vec u$ é o eixo unitário da rotação\fn{Para maior detalhamento, ver \eqref{eq:quaternion_unitario_definicao}.}.

% \subsubsection{Controle e o erro de atitude}
% A formulação da atitude relativa define o erro de orientação em malhas de controle. 
% A partir de $q_{rel}$, é extraído o vetor de erro de atitude $\vec e_q$ e ao combiná-lo com o erro de velocidade angular, obtém-se a correção necessário para o controle aplicar e assim permanece em iteração. 
% Esse procedimento permite que o controlador atue em direção ao menor ângulo de rotação, evitando ambiguidades na representação dos quaternions.

